
RLHF-trained language models exhibit systematic calibration inversion---confidently predicting wrong answers---on specific benchmark task clusters. This work investigates this phenomenon through calibration clustering and mechanism verification. Clustering 2,212 tasks from TruthfulQA, MMLU moral\_scenarios, and Anthropic HH-RLHF by calibration patterns yields a silhouette score of 0.6016 with optimal k=2, revealing non-random failure structure. We trace this pattern to reward signal conflation: annotators rate correctness and user-modeling tasks identically (rate\_diff = 0.001), reward models show similar confidence across task types (mean\_diff = 0.018), and model hidden states fail to separate them (separation = 0.024). Four of five mechanism hypotheses pass validation. However, keyword-based bidirectional feature detection achieves only 2.3\% prevalence with correlation r = -0.027, leaving the feature-cluster link unverified. These findings establish calibration clustering as a diagnostic tool for RLHF behavioral analysis and identify annotator conflation as a plausible mechanism root cause, while highlighting semantic detection as necessary future work.
